
\subsubsection{1.1 Motivation}

Training data deduplication is widely applied in foundation model pretraining. Lee et al. (2022) established that deduplication reduces memorization and improves aggregate benchmark performance. Subsequent corpus construction efforts --- including C4, Dolma, and DCLM --- adopted deduplication as standard practice. The Pythia model suite \citep{Biderman2023Pythia} provides a controlled comparison: two model families trained under otherwise identical conditions (same architecture, optimizer, context length, batch size), differing only in whether the training corpus was deduplicated (Pile versus dedup-Pile).

When Pythia Pile and dedup-Pile models are compared at token-count-matched checkpoints, their benchmark accuracy profiles diverge in a non-uniform way. MMLU accuracy is significantly lower in dedup-Pile models, while HellaSwag and ARC-Challenge scores are higher. This divergence requires explanation. The standard framing --- that deduplication uniformly improves or degrades model quality --- predicts a uniform shift across benchmarks, which is not observed.

\subsubsection{1.2 Research Questions}

This work addresses five specific questions:

\begin{itemize}
\item \textbf{RQ1 (Existence):} Does deduplication produce a statistically significant per-benchmark accuracy differential at token-count-matched checkpoints?
\item \textbf{RQ2 (Correlation):} Does the per-benchmark accuracy differential correlate positively with estimated n-gram contamination levels?
\item \textbf{RQ3 (Methodology):} Does token-count matching recover a stronger contamination signal than step-matching?
\item \textbf{RQ4 (Documents):} Do documents removed by deduplication show higher n-gram overlap with benchmark test content than retained documents?
\item \textbf{RQ5 (Mechanism):} Do Pile-trained models show higher min-k\% probability scores on benchmark items, consistent with near-memorization?
\end{itemize}

\subsubsection{1.3 Hypothesis}

The contamination-correction hypothesis states: if deduplication selectively removes training documents that n-gram-overlap with benchmark test sets, then models trained on the deduplicated corpus should show accuracy reductions proportional to each benchmark's contamination level in the original corpus. Benchmarks more contaminated in the original corpus should show larger accuracy reductions (or smaller increases) after deduplication.

\subsubsection{1.4 Contributions}

\begin{enumerate}
\item \textbf{Empirical --- Contamination-Correlated Benchmark Signature:} Per-benchmark accuracy differentials (dedup-Pile minus Pile) across 16 observations (4 benchmarks x 4 model sizes) correlate positively with estimated 13-gram contamination rates (Pearson r = 0.632, p = 0.0086; Spearman rho = 0.618, p = 0.0107; bootstrap 95\% CI = [0.297, 0.858]). MMLU shows Bonferroni-corrected significant accuracy reduction (t = -5.574, p = 0.0114).
\end{enumerate}

\begin{enumerate}
\item \textbf{Methodological --- Token-Count Matching:} Token-count matching (not step-matching) is the correct confound-control methodology for Pile/dedup-Pile comparisons. Step-matching introduces a volume confound that reduces measured contamination signal by Delta-r = 0.093 and introduces a uniform accuracy bias of approximately -0.004 favoring Pile models in a non-contamination-related manner.
\end{enumerate}

\begin{enumerate}
\item \textbf{Mechanistic (Partial) --- N-gram Overlap of Removed Documents:} A proof-of-concept pipeline confirms that documents removed by deduplication show higher n-gram overlap with benchmark test content than retained documents on 2 of 4 benchmarks (Mann-Whitney U, p < 0.0125 per benchmark, dry-run validation with n = 200 per group, Spearman rho = 1.0 on rank ordering). Full corpus results are pending.
\end{enumerate}

\begin{enumerate}
\item \textbf{Negative Finding --- Min-k\% Direction Reversal at 1B Scale:} The min-k\% membership inference metric applied at Pythia-1B shows the opposite direction from prediction: dedup-Pile models score higher on min-k\% across all four benchmarks (no significant result, 0/4). This does not invalidate the accuracy correlation result but means the near-memorization mechanism is unconfirmed at the model scales tested.
\end{enumerate}

